\subsection*{同一份键到了两个数据库}

先保留 E1 有效、平台余额 10000、各库为空的共同初态。服务原来只有购买库 A。S 的四分量唯一约束把两个并发来件放在同一个仲裁位置，胜者建立 P1/C1/T1，后来的同参来件恢复这份关联。现在为了容纳更多账户，增加购买库 B，并把账户分配到不同分片。分片是保存一部分业务数据及其身份约束的存储单元，A 与 B 各自拥有自己的事务和唯一索引。

假设旧路由把账户 1 送到 A，新路由把它送到 B。一次请求沿缓存中的旧规则到 A，重试沿新规则到 B。两库都查不到 S，各自的唯一约束也都允许首次插入。A 建立 \code{C1/pay-C1}，B 建立 \code{C2/pay-C2}。若平台按各自键正常执行，Alice 的余额依次成为 7000、4000。两个唯一索引都完成了它们各自范围内的仲裁。

\CFourteenFigure{ch14-original01}{路由分歧把同一 S 送入两个独立唯一索引，两份本地首次受理可以形成两项平台付款}{fig:ch14-two-local}

因此需要先回答：S 的首次受理由哪个仲裁域决定？仲裁域是所有可能争夺同一身份的受理者必须共同遵守的决定范围。单库时，它由同一唯一约束及事务提供。分片以后，可以让一个账户的全部购买身份固定归属于同一逻辑分片，也可以另设覆盖完整 S 的全局受理服务。所有能够建立新 P/C 的入口都要进入所选的决定范围。

本章采用按账户固定逻辑归属的方案。账户 1 属于逻辑分片 g7。g7 是操作归属，A、B 是一段时期内承载 g7 的具体存储位置。路由目录保存
\[
g7\longrightarrow(\text{地域 A},\ epoch=8,\ \text{模式 open}).
\]
epoch 是该逻辑分片写入资格的代次，由独立的权威目录从全分片共用的来源递增分配。它区别于 E1 的会话登记代次、T1 的领取代次和数据库的恢复时间线，每一项都比较自己所控制的对象。

把 g7 从 A 搬到 B 时，S 的归属继续是 g7。按地域名重新计算 S 的归属，会使同一操作在切换后落入另一个首次受理空间。修改分片数量也有同样的问题：若直接从 \code{hash(account\_id) mod 4} 改成模 5，部分旧账户会映射到新位置。目录需要保存这批账户的迁移状态，并让新旧路由最终到达同一份身份决定。

\subsection*{重定向以后，接收方依据什么写入}

路由缓存可以很旧，写入资格要在真正受理的位置核验。设目录已经宣布 g7 进入第 9 代，B 是新承载者。请求携带的旧路由提示只能帮助定位。B 从权威资格读取当前归属，再在受保护的入口检查自己的代次和模式。A 收到旧来件时返回重定向或暂时不可用，原 S、原输入和已保存的平台键随恢复路径传递。

第八章用 fencing token 排除旧领取的覆盖写：资源保存已经接受的最大代次 \code{high\_water}，在自己的原子写入边界里比较来件代次、推进水位并修改内容。新持有者取得更大数字以后，还要把水位安装到资源并取得确认，接管完成才包含资源已经排除旧写这一事实。把这项机制用于 g7 时，受保护的资源就是购买库的写入入口。

若购买库有一行 \code{g7,accepted\_epoch=9}，本地写事务先锁住它，比较路由代次，再在同一事务中写 P/C。路由第 8 代事务与把水位推进到 9 的事务通过这把锁排序。已经取得锁的路由第 8 代事务先提交或回滚，推进者等待它结束。水位 9 提交以后，新的路由第 8 代写入在资源处被拒绝。

\CFourteenFigure{ch14-fence-order}{资源行锁把旧事务与水位推进排成次序，第 9 代水位提交以后拒绝旧代写入}{fig:ch14-fence-order}

这说明了这个购买库怎样拒绝旧写。A、B 若各自有一行独立的水位 9，却各自接收同一个 S，两行检查照样可以同时通过。权威目录还要保证同一 g7 在同一受理阶段只授予一个承载者，迁移过程要把已有身份交给后继者。代次比较与身份仲裁分别解决资格和历史的连续性。

可以把一次受理的允许条件写得更准确。设 $\operatorname{owner}(g)$ 是权威目录中的承载者，$e(g)$ 是当前代次，$\operatorname{mode}(g)$ 是准入模式。先只讨论常规开放的入口，r 在以下条件全部成立时才可创建新绑定：
\[
\begin{split}
&\operatorname{owner}(g)=r\ \land\ \operatorname{epoch}(r)=e(g)\\
&{}\land\ \operatorname{mode}(g)=\mathrm{open}\ \land\ \operatorname{eligible}(S).
\end{split}
\]
$\operatorname{eligible}(S)$ 包含账户与动作授权、E1 的活动登记、固定截止和完整 S 的仲裁规则，也包含输入解释与报价核验。切换地域以后，旧会话保留原截止。新地域找不到 E1 时返回过期或恢复中的相应状态，登记器只为独立的新意图分配从未使用的会话。

\subsection*{一次有交接边界的正常迁移}

先讨论历史完全可取得的计划迁移。g7 在 A 的第 8 代开放，S1 已绑定 C1，另有 S2 的受理事务正在执行。迁移者先在权威入口把 g7 改为冻结，阻止新的首次受理，再撤销 A 的本地写资格并等待已经取得资格的事务结束。S2 若提交，它的绑定进入交接范围。若回滚，它没有已提交购买，后继入口以后按同一规则处理。

接着确定交接边界 b。b 表示这批必须带走的权威决定已经完整保存到哪一个位置。若数据库复制负责迁移，边界要对应已确认提交及其恢复位置，接班者重放完这一范围。若业务使用独立有序日志，边界就是它的完整前缀。本章用 $\Gamma$ 表示会话资格和登记进度。迁移者核对业务数据、S 绑定、C/T、原结果、$\Gamma$ 与编号分配进度在同一交接范围内，再让目录授予 B 第 9 代资格。第~\ref{section:14-3}~节会建立一份独立的受理日志作为具体恢复来源。

\CFourteenFigure{ch14-original02}{计划迁移先封住首次受理，收齐第 8 代的在途提交，把完整边界交给 B 后授予第 9 代}{fig:ch14-planned}

B 开放以后，旧路由的 S1 到 A，A 可以把来件引向 B，或给出可重试的迁移响应。B 恢复原 C1。S2 若已在 b 以内提交，也恢复原绑定。尚未受理的独立 S3 在 B 第一次获得资格，才分配新的编号。切换时的短暂停顿，把两次可能首次受理的时间段接成了一个有序边界。

这条迁移的安全条件可以逐项归纳。空系统中每个 S 的绑定数为零。开放阶段只允许当前承载者在唯一约束处从零变为一。冻结阶段禁止新增，允许已排定事务结束并把结果纳入 b。交接复制已有绑定，绑定数保持。后继开放时先读完整交接状态，同 S 进入恢复或冲突分支，只有尚无绑定且具有资格的 S 可以增加一份。归纳成立依赖承载者授予不重叠，以及交接范围确实包含已经确认的事实。

计划迁移最终完成，需要目录能够提交冻结和授予决定，旧事务能够结束，数据通道能够传完边界，B 能够恢复并接受写入。旧存储永久失联时，迁移者可能无法直接取得全部历史。此时选择哪份状态接班，会进入下一节的恢复问题。
\FloatBarrier
